../cictikz/src/cictikz/data/tex/cictikz_preamble.tex

% Switching probability example: two NANDs into a NOR is an AND of all
% four inputs, and the output activity factor follows from the input
% probabilities.

\begin{circuitikz}[thick]

  % the gates
  \node[nand port, color=blue] (n1) at (1.6,0) {};
  \node[nand port, color=blue] (n2) at (1.6,-1.7) {};
  \node[nor port,  color=blue] (n3) at (4.6,-0.85) {};

  \draw[blue] (n1.in 1) -- ++(-0.6,0) node[red, anchor=east] {A};
  \draw[blue] (n1.in 2) -- ++(-0.6,0) node[red, anchor=east] {B};
  \draw[blue] (n2.in 1) -- ++(-0.6,0) node[red, anchor=east] {C};
  \draw[blue] (n2.in 2) -- ++(-0.6,0) node[red, anchor=east] {D};

  \draw[blue] (n1.out) -- ++(0.4,0) |- (n3.in 1);
  \draw[blue] (n2.out) -- ++(0.4,0) |- (n3.in 2);
  \draw[blue] (n3.out) -- ++(0.5,0) node[red, anchor=west] {Y};
  \node[red, anchor=south west] at (2.75,0.05) {X};
  \node[red, anchor=south west] at (2.75,-1.65) {Z};

  % boolean algebra (De Morgan)
  \node[armygreen, anchor=west] at (0,-3.4)
    {$\overline{\overline{AB} + \overline{CD}}
      \;\overset{\mathrm{DM}}{=}\;
      \overline{\overline{AB}} \cdot \overline{\overline{CD}} = ABCD$};
  \node[armygreen, anchor=west] at (0,-4.4)
    {$\overline{D + E} = \overline{D} \cdot \overline{E}$};
  \node[armygreen, anchor=west] at (0,-5.2)
    {$\overline{DE} = \overline{D} + \overline{E}$};

  % probabilities
  \node[armygreen, anchor=west] at (0,-6.3)
    {$P_A = P_B = P_C = P_D = 0.5 = P$};
  \node[armygreen, anchor=west] at (0,-7.3)
    {$P_Y$?};
  \node[armygreen, anchor=west] at (0.5,-8.2)
    {$P_X = 1 - P \cdot P = 1 - \tfrac{1}{4} = \tfrac{3}{4}$};
  \node[armygreen, anchor=west] at (0.5,-9.2)
    {$P_Z = 1 - P \cdot P = 1 - \tfrac{1}{4} = \tfrac{3}{4}$};
  \node[armygreen, anchor=west] at (0.5,-10.2)
    {$\overline{P_X} = \tfrac{1}{4}$};
  \node[armygreen, anchor=west] at (0.5,-11.0)
    {$\overline{P_Z} = \tfrac{1}{4}$};
  \node[armygreen, anchor=west] at (0.5,-12.0)
    {$P_Y = \tfrac{1}{4} \cdot \tfrac{1}{4} = \tfrac{1}{16}$};
  \node[armygreen, anchor=west] at (0.5,-13.1)
    {$\alpha = \tfrac{1}{16}\left(1 - \tfrac{1}{16}\right)
      = \tfrac{15}{16} \cdot \tfrac{1}{16}$};

\end{circuitikz}
\end{document}
